\documentclass[aspectratio=169,usenames,dvipsnames]{beamer}
\beamertemplatenavigationsymbolsempty
% \documentclass{beamer}

% UNCOMMENT FOR HANDOUTS
% \usepackage{pgfpages}
% \pgfpagesuselayout{2 on 1}[landscape,letterpaper,border shrink=5mm]
% \setbeameroption{show notes}

\mode<presentation>
\usetheme{Boadilla}
\usecolortheme{spruce}
\usepackage{../../LizStyle,ulem}
\usepackage{color}
\usepackage[color,matrix,arrow]{xy}
% \input xy
% \xyoption{all}
\usepackage{tikz}
\usepackage{tikz-cd}
\tikzset{commutative diagrams/diagrams={ampersand replacement=\&}}

\AtBeginSection{\frame{\sectionpage}}


% ==========Command for putting comments and removing them=======
% Visible:
\newcommand{\InstructorNotes}[1]{\textit{\color{orange}#1}}
\newcommand{\InstructorList}[1]{
\begin{itemize}
	\color{orange}
	#1
\end{itemize}
}
% Removed:
% \renewcommand{\InstructorNotes}[1]{}
% \renewcommand{\InstructorList}[1]{}


\newcommand{\bvfill}{\vskip0pt plus 1filll}

\usepackage{multimedia}
\usepackage{graphbox}

% \usepackage{algpseudocode}


\newcommand{\Reeb}{\textbf{Reeb}}
\newcommand{\Rtopc}{\textbf{Rtopc}}
% \newcommand{\RR}{\mathcal{R}}

\newcommand{\Set}{\textbf{Set}}
\newcommand{\Vect}{\textbf{Vect}}
\newcommand{\Top}{\textbf{Top}}
\newcommand{\Int}{\textbf{Int}}


\graphicspath{ {Figures/} {../Figures/}{../../Figures/} }



\begin{document}
%\pdfoptionpdfminorversion 5
%\logo{\includegraphics[height=1.5cm]{dukeshield.pdf}}
\title[]{Ch 7.1-7.4: Polynomial regression, Basis functions, Cubic splines}
\subtitle{Lecture 18 - CMSE 381}
\date{Mon, Mar 21, 2022}



\institute[MSU-CMSE]{Michigan State University\\
			::\\
          Dept of Computational Mathematics, Science \& Engineering}
\author[Dr.~Munch]{Prof.~Elizabeth Munch}

\begin{frame}
 \titlepage


\end{frame}
%--------------------------
\begin{frame}[c]{Announcements}

\begin{columns}
\column{.45\textwidth}
\centering

	\textbf{Last time:}
	\begin{itemize}
		\item Lab - PCA/Ridge/Lasso
	\end{itemize}

	\vspace{.2in}

	\textbf{This lecture:}
	\begin{itemize}
	\item 7.1 Polynomial regression
	\item 7.2 Step functions
	\item 7.3 Basis functions
	\item 7.4 Regression Splines
\end{itemize}


\column{.45\textwidth}
\centering

\textbf{Announcements:}
	\begin{itemize}
		\item Mon 3/20: Ch 7.1-7.4
		\item Weds 3/22: Ch 7.5, maybe 7.6?
		\item Friday lab
		\item Friday Homework \#7 due

	\end{itemize}
\end{columns}

\end{frame}
%--- Next Frame ---%



%=================
\section{Last time}
%=================



\begin{frame}[t]{High-Dimensional Data}

\begin{columns}
\column{.45\textwidth}
\centering
\textbf{Low-Dimensions}

$n \gg p$

\footnotesize
\begin{itemize}
\item Low here means $p$ is low, or at least small relative to $n$
\item Can do all the stuff we've talked about so far
\end{itemize}
\column{.45\textwidth}
\centering
\textbf{High-Dimensions}

$n \ll p$

\footnotesize
\begin{itemize}
	\item Issues show up even if $p$ is close to or slightly smaller than $n$
	\item Classical approaches not appropriate since lots of overfitting
\end{itemize}



\end{columns}
\bvfill
\centering
\textbf{What to do about it?}

Be less flexible....
\bvfill
\end{frame}
%--- Next Frame ---%


%--- Next Frame ---%
%--- Next Frame ---%
\begin{frame}[t]{Key points }

\begin{columns}
\column{.45\textwidth}
\centering

\begin{itemize}
  \item regularization
or shrinkage plays a key role in high-dimensional problems, \item  appropriate
tuning parameter selection is crucial for good predictive performance, and
\item  the test error tends to increase as the dimensionality of the problem increases, unless the additional
features are truly associated with the response.
\end{itemize}

\column{.45\textwidth}
\centering
\begin{itemize}
	\item Curse of dimensionality
  \item Report results on an independent test set, or cross-validation errors.
\end{itemize}
\end{columns}

\end{frame}
%--- Next Frame ---%


%==========
\section{Polynomial Regression}
%==========

\begin{frame}[c]{Polynomial regression}

\begin{columns}
\column{.45\textwidth}
\centering
Replace linear model
\begin{equation*}
	y_i = \beta_0 + \beta_1x_1 + \e_i
\end{equation*}
with
\begin{equation*}
	y_i = \beta_0 + \beta_1x_1 + \beta_2x_i^2 + \cdots + \beta_dx_i^d + \e_i
\end{equation*}

\column{.45\textwidth}
\centering
\InstructorList{
  \item Can learn with linear models by passing in predictors $x_i^\ell$
	\item Tend to not go higher than degree 3 or 4 because makes overly flexible
}

\end{columns}

\end{frame}
%--- Next Frame ---%

\begin{frame}[c]{Example with wage data}

\begin{columns}
\column{.45\textwidth}
\centering
\includegraphics[width = \textwidth, align = c]{Fig7-1a}

\column{.45\textwidth}
\centering
\InstructorList{
  \item Plot of wage vs age for men in central Atlantic region of the US
	\item Dark line is degree 4 polynomial
	\item Have variance for each coefficient
	\item Assume you have a model $\hat f(x_0) = \hat \beta_0 + \hat \beta_1 x_0 + \cdots + \hat \beta_d x_0^d$
	\item Can use that to estimate the pointwise variance $\Var(\hat f(x_0))$
	\item Draw 2 std deviations away from the line, 95\% confidence interval
}

\end{columns}

\end{frame}
%--- Next Frame ---%
\begin{frame}[c]{Classification version}

\begin{equation*}
	\Pr(y_i>250 \mid x_i) =
	\frac{\exp(\beta_0 + \beta_1x_i + \cdots + \beta_dx_i^d)}
	     {1+\exp(\beta_0 + \beta_1x_i + \cdots + \beta_dx_i^d)}
\end{equation*}
\begin{columns}
\column{.35\textwidth}
\centering
\includegraphics[width = \textwidth, align = c]{Fig7-1b}
\column{.45\textwidth}
\centering

\footnotesize
\InstructorList{
  \item Note in previous fig that there is a distinct subcluster of high earners making more than \$250K
	\item Build a logistic regression model as above
	\item Note that the 95\% confidence interval gets very wide on the right side
	\item Large sample size $n=3,000$ but small number (79) of high earners makes for high variance in coeffs and wide confidence intervals

}

\end{columns}

\end{frame}
%--- Next Frame ---%
%==========
\section{Step function}
%==========



\begin{frame}[t]{Step functions}
\begin{equation*}
	\mathrm{I}(X < c) \qquad \mathrm{I}(c_1 \leq X < c_2) \qquad \mathrm{I} (c \leq X)
\end{equation*}
\begin{columns}
\column{.45\textwidth}
\centering
\InstructorList{
 \item Draw each of the functions above
}

\column{.45\textwidth}
\centering

\end{columns}

\end{frame}
%--- Next Frame ---%
\begin{frame}[c]{More on step function setup}

\begin{columns}
\column{.45\textwidth}
\centering
\includegraphics[width = \textwidth, align = c]{Eq7-4}
\column{.45\textwidth}
\centering
\InstructorList{
 \item Choose values $c_1,\cdots, c_K$
  \item Allow to learn models that don't have global structure
	\item Use indicator functions to break up the range of $X$ into bins, then we can fit a new constant in each bin.
	\item these are sometimes also called dummy variables
	\item Note that $\sum_j C_j(X) = 1$ because $X$ is in only one interval
}

\end{columns}

\end{frame}
%--- Next Frame ---%
\begin{frame}[c]{Step function: Learned model}

\begin{equation*}
	y_i = \beta_0 + \beta_1C_1(x_i) + \beta_2 C_2(x_i) + \cdots + \beta_K C_K(x_i) + \e_i
\end{equation*}
\begin{columns}
\column{.45\textwidth}
\centering
\InstructorList{
	\item Note above we don't learn a coeff for $C_0$ since like qualitative variable version, we cna figure out hte value from the others.
	\item If $X <c_1$, all predictors are 0 so $\beta_0$ is mean value of $Y$ for $X < c_1$
	\item Then response for $X \in [c_j,c_{j+1})$ is $\beta_0 + \beta_j$, so $\beta_j$ is the avg increase in response for $X$ in the interval relative to $X<c_1$
}


\column{.45\textwidth}
\centering

\end{columns}

\end{frame}
%--- Next Frame ---%

\begin{frame}[c]{Step function example}

\begin{columns}
\column{.45\textwidth}
\centering
\includegraphics[width = \textwidth, align = c]{Fig7-2a}

\column{.45\textwidth}
\centering
\InstructorList{
  \item Learned peak where the high earners are

}

\end{columns}

\end{frame}
%--- Next Frame ---%
\begin{frame}[c]{Step function classification example}
\begin{equation*}
	\Pr(y_i>250 \mid x_i) =
	\frac{\exp(\beta_0 + \beta_1C_1(x_i) + \beta_2 C_2(x_i) + \cdots + \beta_K C_K(x_i))}
	     {1+\exp(\beta_0 + \beta_1C_1(x_i) + \beta_2 C_2(x_i) + \cdots + \beta_K C_K(x_i))}
\end{equation*}
\begin{columns}
\column{.35\textwidth}
\centering
\includegraphics[width = \textwidth, align = c]{Fig7-2b}
\column{.45\textwidth}
\centering
\InstructorList{
  \item Again, learns that middle bit
	\item Still big confidence interval on the right, but likely because there's just less data over there
}

\end{columns}

\end{frame}
%--- Next Frame ---%
\begin{frame}[c]{A few more comments on step functions}

\begin{columns}
\column{.45\textwidth}
\centering
\InstructorList{
  \item Gives the chance to break up the domain, avoid forcing global structure
	\item Need to make decisions about the $c_i$.  A bit arbitrary unless your data has natural breakpoints.
	\item Popular in biostats and epidemiology
}

\column{.45\textwidth}
\centering

\end{columns}

\end{frame}
%--- Next Frame ---%



\begin{frame}[t]{Basis Functions Setup}
Polynomial and piecewise-constant regression models are special cases of a \textit{basis function} approach.

\begin{equation*}
	y_i = \beta_0 + \beta_1b_1(x_i) + \beta_2 b_2(x_i) + \cdots + \beta_K b_K(x_i) + \e_i
\end{equation*}

\begin{columns}
\column{.45\textwidth}
\centering
\InstructorList{
  \item Pick a collection of basis functions $b_1(X),\cdots,b_K(X)$
	\item Explain the $b_i$'s for polynomial and stepwise functions
	\item Use least squares to figure out the constants
	\item Get to still use the inference tools from Ch 3 like standard errors and F-stats

}

\column{.45\textwidth}
\centering

\InstructorList{
	\item Lots of possible options for these
	\InstructorList{
	  \item Some examples are wavelets or Fourier series
		\item Next section is we'll talk about regression splines
	}
}


\end{columns}

\end{frame}
%--- Next Frame ---%
%==========
\section{Regression Splines}
%==========


\begin{frame}[c]{Piecewise polynomials}

\begin{columns}
\column{.45\textwidth}
\centering
\begin{itemize}
	\item Fit a polynomial regression
\begin{equation*}
	y_i = \beta_0 + \beta_1x_1 + \beta_2x_i^2 + \cdots + \beta_dx_i^d + \e_i
\end{equation*}
\item Let the $\beta_i$'s be different at different locations of the range.
\end{itemize}
\column{.45\textwidth}
\centering
\InstructorList{
  \item Points where the coefficients change are called "knots"
	\item A piecewise cubic with no knots is just standard cubic polynomial
	\item More knots gives more flexible polynomial
	\item Piecewise constant of previous section is degree 0 version of this
}

\end{columns}

\end{frame}
%--- Next Frame ---%

\begin{frame}[c]{Example of piecewise polynomial}

\begin{columns}
\column{.45\textwidth}
\centering
\includegraphics[width = \textwidth, align = c]{Fig7-3a}
\column{.45\textwidth}
\centering

Example:
\includegraphics[width = \textwidth, align = c]{Eq7-8b.png}
\end{columns}

\InstructorList{
\item Uses 8 degrees of freedom to fit the model
  \item problem: we have a discontinuous function.

}


\end{frame}
%--- Next Frame ---%

\begin{frame}[c]{The fix}

\begin{columns}
\column{.45\textwidth}
\centering
\begin{itemize}
	\item Fit piecewise polynomial
	\item Require continuity at knots
\end{itemize}
\column{.45\textwidth}
\centering
\includegraphics[width = \textwidth, align = c]{Fig7-3b}
\end{columns}
\InstructorNotes{Looks better, but pointy join is odd}
\end{frame}
%--- Next Frame ---%

\begin{frame}[t]{The better fix: Cubic splines}

\begin{columns}
\column{.45\textwidth}
\centering

\begin{itemize}
	\item Fit piecewise polynomial
	\item Require continuity at knots
	\item Require the first and second derivatives to be continuous at knots
\end{itemize}
\InstructorList{
  \item Now we have smooth joins
	\item This is called a "cubic spline"
	\item 8 original degrees of freedom, plus three constraints (continuous, 1st deriv continuous, 2nd deriv continuous) leaves 5 degrees of freedom
	\item cubic spline with $K$ knots uses $4+K$ degrees of freedom
}
\column{.45\textwidth}
\centering
\includegraphics[width = \textwidth, align = c]{Fig7-3c}
\end{columns}

\end{frame}
%--- Next Frame ---%
\begin{frame}[c]{Generalized splines}

\begin{columns}
\column{.45\textwidth}
\centering
\InstructorList{
  \item Can do this for any degree $d$ polynomial
	\item Require continuity up to degree $d-1$ at each knot
	\item Linear then is just requiring continuity
}

\column{.45\textwidth}
\centering
\includegraphics[width = \textwidth, align = c]{Fig7-3d}
\end{columns}

\end{frame}
%--- Next Frame ---%

\begin{frame}[c]{Spline basis representation}

\begin{columns}
\column{.65\textwidth}
\centering
Want to pick $b_i$ so that we represent a cubic spline with $K$ knots as
\begin{equation*}
	y_i = \beta_0 + \beta_1b_1(x_i) + \beta_2 b_2(x_i) + \cdots + \beta_{K+3} b_{K+3}(x_i) + \e_i
\end{equation*}
\end{columns}

\end{frame}
%--- Next Frame ---%


\begin{frame}[t]{Truncated power basis function}

\begin{columns}
\column{.45\textwidth}
\centering

\begin{equation*}
h(x,z) = (x-z)^3_+ =
\begin{cases}
	(x-z)^3 & \text{if } x>z \\
	0 & else
\end{cases}
\end{equation*}
Desmos link:
\url{https://www.desmos.com/calculator/esucuulbgj}
\column{.45\textwidth}
\centering
\InstructorList{
  \item Book uses $z=\xi$ notation, but i refuse
	\item Draw $h(x,2)$
}

\end{columns}

\end{frame}
%--- Next Frame ---%
\begin{frame}[t]{The basis for cubic splines }

\begin{columns}
\column{.45\textwidth}
\centering
Given knots at $z_1,\cdots,z_K$
\begin{itemize}
	\item $X$
	\item $X^2$
	\item $X^3$
	\item $h(X,z_1)$
	\item $h(X,z_2)$
	\item $\vdots$
	\item $h(X,z_K)$
\end{itemize}
\column{.45\textwidth}
\centering


\InstructorList{
  \item Adding $h(X,z_i)$ leads to discontinuity at most at the third derivative at $z_i$
	\item So use the $3+K$ functions listed here, result is $K+4$ coefficients once you get the intercept.
	\item Matches with our $K+4$ degrees of freedom from earlier
}

\end{columns}

\end{frame}
%--- Next Frame ---%

\begin{frame}[t]{Notes on cubic splines }

\begin{columns}
\column{.45\textwidth}
\centering

\InstructorList{
  \item Human eyes can't detect the discontinuities so often cubic splines are a good way to go
	\item Still have problem of high variance at the outer range of predictors (big or small $X$)
	\item \textit{Natural spline:} regression spline plus boundary constraints where function is required to be linear at the boundary
}

\column{.45\textwidth}
\centering

\includegraphics[width = \textwidth]{Fig7-4}
\end{columns}

\end{frame}
%--- Next Frame ---%
\begin{frame}[t]{Where to put the knots?}

\begin{columns}

\column{.65\textwidth}
\centering

\includegraphics[width = \textwidth]{Fig7-5}
\column{.25\textwidth}
\centering

\InstructorList{
  \item Plan A: Uniform placement
	\item Plan B: Data driven. Example here puts it at quantiles of the data
}
\end{columns}

\end{frame}
%--- Next Frame ---%
\begin{frame}[t]{How many knots to use?}
{When in doubt, Cross-Validate.}
\begin{columns}
\column{.45\textwidth}
\centering

\InstructorList{
  \item Do CV ($k$-fold probably) on data for each number of knots
	\item Figure at right is the MSE resulting from 10-fold CV
	\item Note that clearly linear is not enough, after that 3-4 degrees of freedom for the cubic spline are enough
}

\column{.45\textwidth}
\includegraphics[width = \textwidth]{Fig7-6b}
\centering

\end{columns}
\end{frame}
%--- Next Frame ---%
\begin{frame}[t]{Cubic splines vs Polynomial Regression}

\begin{columns}
\column{.45\textwidth}
\centering
\includegraphics[width = \textwidth]{Fig7-7}
\column{.45\textwidth}
\centering

\InstructorList{
  \item Left: natural cubic spline with 15 degrees of freedom vs degree 15 polynomial
	\item Polynomial flexibility means ugly boundary behavior
	\item Splines let you add flexibility as more knots rather than more degrees, so allows for flexibility where the function is changing more rapidly
}

\end{columns}

\end{frame}
%--- Next Frame ---%


\begin{frame}[t]{Next time}


\begin{columns}
\column{.45\textwidth}
\centering

\textbf{Next time:}

\begin{itemize}
	\item 7.5, Maybe 7.7
\end{itemize}

\column{.45\textwidth}
\centering

\textbf{This week:}
\begin{itemize}

	\item Mon 3/14: 7.1-7.4
	\item Weds 3/16: 7.5, maybe 7.7?
	\item Fri 3/18: Lab: 7.8
	\begin{itemize}
		\item HW \#7 Due Friday
		\item Covers PCA/Ridge/Lasso Lab. Also content from today and Wednesday.
	\end{itemize}
\end{itemize}
	\end{columns}

\end{frame}
%--- Next Frame ---%


\end{document}
